% sections/heterogeneity.tex: Section 6: exposure at population scale, the term margin at population scale against the within-borrower
% design, the geography of exposure and its correlates, and exposure against impatience. Results: panel_windows_auto.csv,
% panel_power_splits_auto.csv (T1l), heterogeneity_population.csv (heterogeneity_population.R), geo_correlates_v3.csv (geo_correlates_v3.R),
% identification_person_place_v2.csv (I7), mc_rho_variance*.csv; memo memos/heterogeneity_program_2026-09-16.md.
\section{Heterogeneity at Population Scale}\label{sec:heterogeneity}

The population specification of equation \eqref{eq:main_reg} can be estimated within any group of loans, and the panel records the credit score, income, the lender, the zip code and the origination date of every loan. Its identification is weaker than the within-borrower design's and the experiments': it rests on the assumption of Section \ref{sec:identification}, and its cells were selected by the rule stated there, keep the lender's menu and remove the person, which is the rule under which the population payment response matches the within-borrower one and the published quasi-experiments. The step is worth its assumption because of what it buys. The within-borrower design reports ten score deciles, four income quartiles, five age groups, five periods and twelve states; the population design reports the same objects at every resolution the panel supports, and three results of this section exist only there: the geography of exposure below the state, the shape of exposure across the score distribution at equal power, and what the geography is correlated with. What the step does not buy is a rate by group. Section \ref{sec:termmargin_pop} shows that the population term margin is a selection gradient in the score, and the rule of Section \ref{sec:rate_people}, no rate for a place from the term margin, holds at population scale as well; the rate rests on the within-borrower design. The section reports the payment elasticity by group, which the model interprets as the group's \emph{exposure} to the default margin, the density of its members at the default threshold times the weight of the payment stream (Appendix \ref{appendix:aggregation}); exposure differs across groups by a factor of two, the rate of Section \ref{sec:rate_people} does not, and Section \ref{sec:exposure_not_impatience} states what follows.

\subsection{Exposure Across People and Places}\label{sec:exposure}

Figure \ref{fig:exposure_groups} in Appendix \ref{appendix:figures} reports the elasticity of default within 24 months to the payment by score decile, income quartile, census region and origination cohort, with the within-borrower estimate as the reference line, and Table \ref{tab:het_population} gives the values with the default rate of each group. The elasticity rises from $0.41$ in the bottom income quartile to $0.77$ in the third and $0.86$ in the fourth, and from $0.48$ in the bottom score tercile to $0.71$ in the middle, with the score-decile profile peaking at $0.80$ in the fifth decile and falling to $0.31$ in the ninth, where defaults are rare; across regions it is $0.62$ in the Northeast and $0.41$ to $0.47$ elsewhere; across origination cohorts it is $0.37$ to $0.40$ in every cohort from 2010 onward against $0.59$ before; and it rises with the contract's length, from $0.23$ for contracts of 48 months or less to $0.78$ for 73 months or more. The income gradient is not score composition: within the bottom score tercile exposure rises from $0.44$ in the bottom income quartile to $0.97$ in the top, within the middle tercile from $0.67$ to $0.88$ and within the top from $0.32$ to $0.67$, cells the within-borrower design cannot report. Low-income and low-score borrowers respond \emph{less} to a given proportional change in the payment, which is the pattern of a population whose members are on average further from the threshold rather than concentrated at it.

The finer cuts are power-preserving, and they are what the population step is for. Figure \ref{fig:exposure_geo} in Appendix \ref{appendix:figures} splits the score and income distributions into 46 bins each built to hold an equal number of charge-offs over the life of the loan, 5,300 to 8,000, so that every bin's standard error is $0.06$ to $0.11$, and clusters three-digit zip codes into 91 areas holding 1,400 to 7,700 lifetime charge-offs each. Across the score bins exposure rises from $0.47$ at a mean score of 365 to $0.82$ at 690 and falls to $0.45$ at 795; the within-borrower design's deciles cannot place the peak. Across the 91 areas the elasticity runs from $0.33$ at the tenth percentile to $0.67$ at the ninetieth, a ratio of two, with a loan-weighted standard deviation of $0.125$ against a root mean squared standard error of $0.064$: the dispersion net of sampling error, $0.107$, is as large as the dispersion across income quartiles and larger than across cohorts. Appendix \ref{appendix:traits} puts the place's share of a group's default rate at 7 to 12 percent, from borrowers who move, and the person's share of the variance of default at a quarter. Section \ref{sec:geo_correlates} asks what the geography is correlated with.

\subsection{The Term Margin at Population Scale}\label{sec:termmargin_pop}

Table \ref{tab:het_population} sets the population term margin beside the within-borrower one, group by group: the ratio of the term response to the payment response at 24 months, the rate it maps to through equation \eqref{eq:window} with the group's own profile and exit hazard, the same rate after the selection adjustment of Appendix \ref{appendix:ladder}, which adds to the population term response the within-minus-across difference measured on the multi-loan borrowers of the same group, and the within-borrower rate of Table \ref{tab:rho_within}. Pooled, the population ratio is $0.25$ ($0.09$) against $0.55$ ($0.01$) within borrower, a raw rate of $0.34$ with a set from $0.14$ to $0.77$; the selection difference closes a quarter of the gap, and the adjusted rate, $0.13$ with a set from zero to $0.34$, contains the within-borrower estimate. By group the two designs part company, and the way they part is informative. Across score deciles the population ratio rises monotonically, from $-0.95$ in the bottom decile through zero in the third to $4.45$ in the ninth, while the within-borrower ratio falls from $0.53$ to $0.32$ at the sixth decile and returns to $0.49$ at the top; at the resolution of the 46 equal-power score bins the population ratio rises with the score at a rank correlation of $0.89$ and crosses zero at a score of 605, and across the 46 income bins it rises from $-0.80$ to $0.89$. That is the selection pattern of \citet{hertzberg_screening_2018}, lenders screening risky borrowers with short maturities and safe borrowers choosing long ones, measured at 46 points of the score distribution on seven million loans, where the within-borrower design measures its size at ten. A single selection difference per group does not turn the population term response into a within-borrower one: the adjusted population rate overlaps the within-borrower set in four deciles and in every income quartile and sits at the ceiling in the middle deciles. By age the two designs agree in direction and not in level: both ratios rise with age, the population's from $-0.20$ for ages 18 to 30 to $1.02$ above 60 and the within-borrower's from $0.26$ to $0.65$, so both find the youngest borrowers weighing the far payments least, while the population rate is undefined below age 41, where the term response is negative, and its sets overlap the within-borrower sets in the two oldest groups only. By state (Table \ref{tab:het_geo}, Appendix \ref{appendix:figures}) the ratio's standard errors are $0.09$ to $0.20$ against $0.01$ pooled within borrower, and the two designs' point values are unrelated. Figure \ref{fig:het_agreement} in Appendix \ref{appendix:figures} shows the three rates by decile, quartile and state. The population step delivers exposure with power everywhere and the selection gradient at fine resolution; it does not deliver a rate by group, and the paper does not report one.

\begin{table}[H]
\centering
\caption{The population design beside the within-borrower design, by group}
\label{tab:het_population}
\scriptsize
\resizebox{\textwidth}{!}{\begin{tabular}{llrrrrrrr}
\toprule
 & & & Default & Exposure & Ratio & \multicolumn{2}{c}{Population rate} & Within-borrower \\
\cmidrule(lr){7-8}
Dimension & Group & Loans & rate (\%) & $\varepsilon_m$ (s.e.) & $R$ (s.e.) & raw & selection-adjusted & rate \\
\midrule
All loans &  & 6,971,966 & 1.86 & 0.45 (0.02) & 0.25 (0.09) & 0.34 [0.14, 0.77] & 0.13 [0.00, 0.34] & 0.02 [0.01, 0.04] \\
\addlinespace
Score decile & 1 & 696,601 & 7.79 & 0.46 (0.03) & -0.95 (0.11) & -- & 0.36 [0.17, 0.73] & 0.14 [0.10, 0.18] \\
 & 2 & 702,476 & 4.75 & 0.53 (0.03) & -0.43 (0.08) & -- & 0.08 [0.00, 0.22] & 0.11 [0.08, 0.16] \\
 & 3 & 705,619 & 2.66 & 0.63 (0.03) & 0.04 (0.08) & 1.03 [0.44, $\infty$] & 0.00 [0.00, 0.05] & 0.10 [0.06, 0.14] \\
 & 4 & 703,860 & 1.55 & 0.71 (0.03) & 0.35 (0.09) & 0.20 [0.03, 0.46] & 0.00 [0.00, 0.02] & 0.11 [0.07, 0.16] \\
 & 5 & 701,698 & 0.84 & 0.80 (0.05) & 0.61 (0.09) & 0.00 [0.00, 0.10] & 0.00 [0.00, 0.00] & 0.17 [0.12, 0.23] \\
 & 6 & 699,576 & 0.45 & 0.72 (0.05) & 0.90 (0.13) & 0.00 [0.00, 0.00] & 0.00 [0.00, 0.14] & 0.22 [0.15, 0.29] \\
 & 7 & 698,951 & 0.24 & 0.69 (0.07) & 1.44 (0.21) & 0.00 [0.00, 0.00] & 0.00 [0.00, 0.05] & 0.12 [0.06, 0.20] \\
 & 8 & 700,541 & 0.12 & 0.49 (0.10) & 2.17 (0.52) & 0.00 [0.00, 0.00] & 0.00 [0.00, 0.48] & 0.10 [0.02, 0.19] \\
 & 9 & 696,650 & 0.07 & 0.31 (0.10) & 4.45 (1.54) & 0.00 [0.00, 0.00] & 0.00 [0.00, 0.48] & 0.05 [0.00, 0.18] \\
\addlinespace
Score tercile & 1 & 2,339,724 & 4.74 & 0.48 (0.02) & -0.48 (0.07) & -- & 0.07 [0.00, 0.19] & 0.11 [0.09, 0.13] \\
 & 2 & 2,336,136 & 0.72 & 0.71 (0.03) & 0.71 (0.08) & 0.00 [0.00, 0.00] & 0.00 [0.00, 0.00] & 0.15 [0.12, 0.18] \\
 & 3 & 2,296,106 & 0.09 & 0.47 (0.05) & 2.73 (0.38) & 0.00 [0.00, 0.00] & 0.00 [0.00, 0.00] & 0.06 [0.01, 0.11] \\
\addlinespace
Income quartile & 1 & 1,732,024 & 4.94 & 0.41 (0.02) & -0.51 (0.08) & -- & 0.18 [0.01, 0.47] & 0.24 [0.15, 0.36] \\
 & 2 & 1,767,854 & 1.78 & 0.61 (0.03) & 0.26 (0.07) & 0.33 [0.16, 0.63] & 0.00 [0.00, 0.00] & 0.06 [0.00, 0.15] \\
 & 3 & 1,768,129 & 0.53 & 0.77 (0.04) & 0.66 (0.07) & 0.00 [0.00, 0.01] & 0.18 [0.00, 0.53] & 0.08 [0.00, 0.21] \\
 & 4 & 1,703,959 & 0.19 & 0.86 (0.05) & 0.77 (0.10) & 0.00 [0.00, 0.00] & -- & 0.26 [0.02, 0.92] \\
\addlinespace
Score tercile $\times$ income quartile & 1 / 1 & 1,187,381 & 6.62 & 0.44 (0.02) & -0.78 (0.09) & -- & 0.23 [0.08, 0.44] & -- \\
 & 1 / 2 & 804,978 & 3.20 & 0.64 (0.03) & -0.13 (0.07) & -- & 0.00 [0.00, 0.08] & -- \\
 & 1 / 3 & 260,739 & 2.02 & 0.83 (0.05) & -0.10 (0.10) & -- & 0.07 [0.00, 0.28] & -- \\
 & 1 / 4 & 86,626 & 1.39 & 0.97 (0.11) & -0.11 (0.20) & -- & 0.15 [0.00, 1.44] & -- \\
 & 2 / 1 & 524,622 & 1.34 & 0.67 (0.04) & 0.47 (0.11) & 0.08 [0.00, 0.31] & 0.08 [0.00, 0.31] & -- \\
 & 2 / 2 & 724,088 & 0.73 & 0.74 (0.05) & 0.80 (0.11) & 0.00 [0.00, 0.00] & 0.00 [0.00, 0.00] & -- \\
 & 2 / 3 & 644,007 & 0.48 & 0.85 (0.06) & 0.78 (0.12) & 0.00 [0.00, 0.00] & 0.00 [0.00, 0.00] & -- \\
 & 2 / 4 & 443,419 & 0.33 & 0.88 (0.08) & 0.83 (0.16) & 0.00 [0.00, 0.02] & 0.00 [0.00, 0.02] & -- \\
 & 3 / 2 & 238,788 & 0.17 & 0.32 (0.15) & 2.53 (1.32) & 0.00 [0.00, $\infty$] & 1.21 [0.00, $\infty$] & -- \\
 & 3 / 3 & 863,383 & 0.11 & 0.39 (0.08) & 3.47 (0.81) & 0.00 [0.00, 0.00] & 0.00 [0.00, 0.00] & -- \\
 & 3 / 4 & 1,173,914 & 0.05 & 0.67 (0.10) & 2.02 (0.39) & 0.00 [0.00, 0.00] & 0.00 [0.00, 0.26] & -- \\
\addlinespace
Age & 18-30 & 1,393,345 & 3.16 & 0.53 (0.02) & -0.20 (0.09) & -- & 0.48 [0.18, $\infty$] & 0.32 [0.21, 0.48] \\
 & 31-40 & 1,512,490 & 2.15 & 0.47 (0.02) & -0.03 (0.10) & -- & 0.43 [0.09, $\infty$] & 0.00 [0.00, 0.10] \\
 & 41-50 & 1,571,232 & 1.53 & 0.42 (0.02) & 0.30 (0.12) & 0.26 [0.04, 0.76] & 0.29 [0.00, $\infty$] & 0.00 [0.00, 0.03] \\
 & 51-60 & 1,339,966 & 1.12 & 0.50 (0.03) & 0.56 (0.10) & 0.01 [0.00, 0.18] & 0.26 [0.00, $\infty$] & 0.00 [0.00, 0.00] \\
 & 61+ & 1,114,436 & 1.18 & 0.55 (0.03) & 1.02 (0.10) & 0.00 [0.00, 0.00] & 0.21 [0.00, $\infty$] & 0.00 [0.00, 0.07] \\
\addlinespace
Census region & Northeast & 1,107,525 & 1.25 & 0.62 (0.04) & 0.59 (0.12) & 0.00 [0.00, 0.20] & 0.00 [0.00, 0.19] & -- \\
 & Midwest & 1,515,086 & 1.43 & 0.44 (0.02) & 0.29 (0.11) & 0.28 [0.06, 0.77] & 0.26 [0.04, 0.76] & -- \\
 & South & 2,794,599 & 2.40 & 0.41 (0.02) & 0.05 (0.08) & 0.92 [0.43, $\infty$] & 0.84 [0.38, $\infty$] & -- \\
 & West & 1,503,243 & 1.77 & 0.47 (0.02) & 0.12 (0.09) & 0.61 [0.28, $\infty$] & 0.58 [0.25, $\infty$] & -- \\
\addlinespace
Origination cohort & 2005-09 & 1,966,929 & 1.25 & 0.59 (0.05) & -0.07 (0.11) & -- & -- & -- \\
 & 2010-12 & 1,206,366 & 1.24 & 0.39 (0.04) & -0.03 (0.16) & -- & -- & -- \\
 & 2013-15 & 1,064,256 & 2.40 & 0.40 (0.04) & 0.30 (0.21) & 0.27 [0.00, $\infty$] & 0.26 [0.00, $\infty$] & -- \\
 & 2016-18 & 1,100,059 & 2.42 & 0.40 (0.04) & 0.58 (0.28) & 0.02 [0.00, 1.01] & 0.01 [0.00, 0.95] & -- \\
 & 2019-21 & 1,038,423 & 1.90 & 0.38 (0.04) & 0.40 (0.21) & 0.17 [0.00, $\infty$] & 0.16 [0.00, $\infty$] & -- \\
 & 2022+ & 595,933 & 3.07 & 0.37 (0.05) & 0.26 (0.21) & 0.37 [0.00, $\infty$] & 0.35 [0.00, $\infty$] & -- \\
\addlinespace
Contract length (months) & $\le$48 & 1,887,241 & 1.61 & 0.23 (0.02) & 0.07 (0.22) & 0.80 [0.05, $\infty$] & 0.70 [0.02, $\infty$] & -- \\
 & 49-60 & 1,708,168 & 1.22 & 0.44 (0.02) & -1.94 (0.69) & -- & -- & -- \\
 & 61-72 & 2,318,901 & 2.31 & 0.60 (0.03) & 1.08 (0.27) & 0.00 [0.00, 0.05] & 0.00 [0.00, 0.04] & -- \\
 & 73+ & 1,057,656 & 2.35 & 0.78 (0.04) & -3.37 (0.38) & -- & -- & -- \\
\bottomrule
\end{tabular}
}
\vspace{0.5em}
\begin{minipage}{0.95\textwidth}
\noindent \small \textit{Notes:} Auto loans in the monthly one-percent panel, charge-off within 24 months, the population design of equation \eqref{eq:main_reg} within each group. Exposure is the payment elasticity $\varepsilon_m$; the ratio $R$ is the term elasticity over the payment elasticity, with the delta-method standard error. The population rate maps $R$ through equation \eqref{eq:window} at $K = 24$ with the group's own marginal-utility profile, exit hazard and mean contract length; the selection-adjusted rate first adds to the term elasticity the within-borrower minus across-borrower difference of Table \ref{tab:within_by_group} for the group (the pooled difference where the group has none); a ratio at or below zero has no rate (--), and a ratio at or above the formula's ceiling gives zero. The within-borrower rate is Table \ref{tab:rho_within}. Brackets are 95 percent sets. The tenth score decile has too few charge-offs within 24 months for the design; the score-tercile-by-income-quartile cells and the age groups (age at first sight) exist only at population scale; the contract-length bins restrict the term to a few menu points and their ratios are not term margins.
\end{minipage}
% replication: Code/type/T1l_windows_panel.R, Code/local/heterogeneity_population.R, Code/exhibits/make_heterogeneity_tables.R; results/heterogeneity_population.csv
\end{table}

\subsection{The Geography of Exposure and Its Correlates}\label{sec:geo_correlates}

Table \ref{tab:geo_correlates} relates the three population objects, exposure, the log default rate and the population rate, to the economy and the vote of the place: adjusted gross income per return, the wage share of income, the shares of returns with unemployment compensation and with the earned income credit, the occupational exposure of employment to artificial intelligence, the mean unemployment rate of 2010 to 2024, and the Republican two-party share of the presidential vote in 2016 and 2020 and its change from 2012. Each cell is the effect of one loan-weighted standard deviation of the covariate, raw and with the place's composition (its mean score, income and contract length) held fixed, across the 43 states with covariates (Panel A) and the 91 zip-code clusters (Panel B). Three results follow. The geography of default is the geography of low income: one standard deviation of the earned-income-credit share raises the log default rate by $0.26$ across states and $0.28$ across clusters, four fifths of its dispersion at either level, and the place's composition absorbs it ($0.06$ and $0.07$ with controls). The geography of exposure is, at the finer level, composition too: across states one standard deviation of income per return raises exposure by $0.07$, two thirds of its cross-state standard deviation, and the relation survives the state's composition ($0.068$, standard error $0.017$); across the 91 clusters the raw gradient is $0.047$ and it is $0.000$ ($0.018$) with the cluster's composition held fixed, so the income gradient in exposure across places is the score, income and contract length of the people who live there, with a residual across states that the clusters do not show. The vote is composition at both levels: places that vote Republican have lower exposure raw, $-0.06$ per standard deviation of the 2020 share across states and across clusters, and the same exposure as everyone else with composition held fixed ($-0.004$ and $-0.024$); the earned-income-credit share, the unemployment-compensation share and the AI-exposure index behave as the vote does. The population rate is related to nothing once composition is held fixed at either level, every coefficient within $1.6$ standard errors of zero, which is the term margin's selection gradient seen through places. Figure \ref{fig:geo_correlates} in Appendix \ref{appendix:figures} shows the two raw gradients and the cloud. Places that vote Republican hold more households near the default threshold because they are poorer, and they weigh far payments as everyone else does; the political map of default is an income map.

\begin{table}[H]
\centering
\caption{What the geography of exposure, default and the population rate is correlated with}
\label{tab:geo_correlates}
\footnotesize
\resizebox{\textwidth}{!}{\begin{tabular}{lcccccc}
\toprule
 & \multicolumn{2}{c}{Exposure $\varepsilon_m$} & \multicolumn{2}{c}{Default rate (log)} & \multicolumn{2}{c}{Population rate} \\
\cmidrule(lr){2-3}\cmidrule(lr){4-5}\cmidrule(lr){6-7}
Covariate (one s.d.) & raw & composition & raw & composition & raw & composition \\
\midrule
\multicolumn{7}{l}{\textit{Panel A: states (43 units; 30 with a rate)}} \\
\multicolumn{7}{l}{\quad \textit{Economic}} \\
\quad AGI per return (log) & 0.072$^{***}$ (0.014) & 0.068$^{***}$ (0.017) & -0.078$^{*}$ (0.046) & 0.038$^{*}$ (0.021) & -0.004 (0.053) & -0.030 (0.110) \\
\quad Wage share of AGI & -0.023$^{**}$ (0.010) & -0.038$^{***}$ (0.006) & -0.063$^{*}$ (0.032) & -0.020$^{**}$ (0.010) & -0.051$^{**}$ (0.024) & -0.015 (0.026) \\
\quad Returns with unemployment compensation & 0.061$^{***}$ (0.018) & 0.001 (0.018) & -0.147$^{**}$ (0.072) & 0.067$^{**}$ (0.031) & -0.120$^{**}$ (0.057) & -0.221 (0.190) \\
\quad EITC share of returns & -0.046$^{***}$ (0.015) & 0.040$^{*}$ (0.021) & 0.259$^{***}$ (0.028) & 0.058 (0.042) & 0.099 (0.081) & -0.091 (0.112) \\
\quad AI exposure of employment & 0.056$^{***}$ (0.017) & 0.013 (0.017) & -0.131$^{***}$ (0.037) & -0.065$^{**}$ (0.030) & -0.041 (0.085) & 0.029 (0.087) \\
\quad Unemployment rate, 2010--2024 & 0.031$^{**}$ (0.012) & 0.011 (0.014) & 0.068 (0.043) & 0.076$^{***}$ (0.013) & -0.052 (0.058) & -0.078 (0.079) \\
\multicolumn{7}{l}{\quad \textit{Political}} \\
\quad Republican share, 2016 & -0.061$^{***}$ (0.014) & -0.012 (0.022) & 0.103$^{**}$ (0.049) & -0.053$^{**}$ (0.025) & 0.069 (0.050) & 0.084 (0.141) \\
\quad Republican share, 2020 & -0.060$^{***}$ (0.014) & -0.004 (0.021) & 0.114$^{**}$ (0.047) & -0.037 (0.026) & 0.058 (0.052) & 0.042 (0.119) \\
\quad Republican share, 2016 minus 2012 & 0.003 (0.018) & 0.011 (0.013) & -0.084 (0.053) & -0.035 (0.029) & -0.079 (0.065) & 0.077 (0.122) \\
\midrule
\multicolumn{7}{l}{\textit{Panel B: equal-default zip3 clusters (91 units; 68 with a rate)}} \\
\multicolumn{7}{l}{\quad \textit{Economic}} \\
\quad AGI per return (log) & 0.047$^{***}$ (0.016) & 0.000 (0.018) & -0.069$^{**}$ (0.031) & 0.000 (0.017) & 0.012 (0.039) & -0.035 (0.050) \\
\quad Wage share of AGI & -0.006 (0.011) & -0.008 (0.010) & -0.023 (0.027) & -0.015 (0.009) & -0.024 (0.029) & -0.025 (0.030) \\
\quad Returns with unemployment compensation & 0.060$^{***}$ (0.014) & 0.028$^{**}$ (0.012) & -0.156$^{***}$ (0.033) & 0.026 (0.020) & -0.068 (0.046) & -0.100 (0.062) \\
\quad EITC share of returns & -0.037$^{**}$ (0.015) & 0.040$^{*}$ (0.022) & 0.275$^{***}$ (0.022) & 0.072$^{***}$ (0.026) & -0.011 (0.036) & -0.095 (0.071) \\
\quad AI exposure of employment & 0.043$^{***}$ (0.014) & -0.005 (0.015) & -0.102$^{***}$ (0.028) & -0.047$^{***}$ (0.018) & 0.007 (0.042) & -0.031 (0.053) \\
\quad Unemployment rate, 2010--2024 & 0.023$^{**}$ (0.011) & 0.005 (0.012) & 0.067$^{*}$ (0.036) & 0.050$^{***}$ (0.013) & 0.020 (0.034) & 0.023 (0.037) \\
\multicolumn{7}{l}{\quad \textit{Political}} \\
\quad Republican share, 2016 & -0.059$^{***}$ (0.012) & -0.026$^{*}$ (0.015) & 0.040 (0.028) & -0.031 (0.020) & 0.005 (0.036) & 0.032 (0.050) \\
\quad Republican share, 2020 & -0.058$^{***}$ (0.013) & -0.024 (0.015) & 0.060$^{**}$ (0.027) & -0.019 (0.021) & 0.003 (0.037) & 0.035 (0.051) \\
\quad Republican share, 2016 minus 2012 & -0.006 (0.014) & -0.005 (0.017) & -0.111$^{**}$ (0.044) & -0.035$^{*}$ (0.018) & -0.105$^{**}$ (0.047) & -0.035 (0.057) \\
\bottomrule
\end{tabular}
}
\vspace{0.5em}
\begin{minipage}{0.95\textwidth}
\noindent \small \textit{Notes:} Loan-weighted regressions across the 43 states with covariates (30 with a population rate; Panel A) and the 91 equal-default clusters of three-digit zip codes (68 with a rate; Panel B) of the unit's exposure, log 24-month charge-off rate and population rate on one covariate at a time; the cell is the standardized coefficient, the effect of one loan-weighted standard deviation of the covariate in the outcome's units, with the HC1 standard error; ``composition'' adds the state's mean credit score, mean income and mean contract length. Covariates: IRS Statistics of Income by zip code for 2022, aggregated to the unit; the Felten--Raj--Seamans occupational AI-exposure index weighted by employment; the mean state unemployment rate over 2010 to 2024; county presidential returns aggregated to the unit. Specifications with census-region effects are in the results file. $^{*}$, $^{**}$, $^{***}$: 10, 5 and 1 percent.
\end{minipage}
% replication: Code/local/geo_correlates_v3.R, Code/exhibits/make_geo_figures.R; results/geo_correlates_v3.csv, geo_correlates_v3_units.csv
\end{table}

\subsection{Exposure, Not Impatience}\label{sec:exposure_not_impatience}

The results of Sections \ref{sec:rate_people} and \ref{sec:exposure} support one statement about the heterogeneity in default, and this subsection states it and says what it rests on. Each month a borrower defaults when the gain from not paying, the payment she keeps, exceeds what she gives up: the penalty, and the value of staying current, which is the stream of the payments she would otherwise make, each weighted by its distance in the manner of Proposition \ref{prop:sensT}. Two objects therefore determine how many members of a group default and how that number moves when the payment changes. The first is the location of the group's members relative to the threshold. The default rate is the mass across the line, and the response to a change in the payment is the density at it. That is exposure. The second is how the group weighs the future against the present, which sets the value of staying current: the weight on each future payment, the discount function times survival times the marginal-utility ratio. That is impatience, and the present-payment premium.

The two have different observable consequences, and the data show differences in the first and not the second. Exposure moves the level of the payment elasticity: more borrowers at the line, a larger response per unit of default. The elasticity is $0.41$ in the bottom income quartile and $0.86$ in the top, and the score profile peaks at the median. The rate, from the within-borrower design, varies non-monotonically with score and income, steepest near the median score and at the lowest incomes, and nowhere rises with the group's default rate (Section \ref{sec:rate_people}); and the one experiment that reports both of its arms by subgroup finds the four-year rate the same above and below the median credit score and the median debt-to-income ratio (Section \ref{sec:quasi}, Table \ref{tab:rho_subgroups}). The level pattern says something specific about the location of each group relative to the threshold. A low-income borrower who defaults is, on average, further past the threshold than a high-income borrower who defaults: more of her defaults are forced by an income shortfall, and fewer are marginal choices that a ten percent lower payment would have reversed. Her group's default rate within 24 months is twenty-six times that of the top quartile and its elasticity is half as large. Both facts come from the distribution of income and buffers relative to the payment and the penalty, the location of the group around the line, and neither requires that she weigh the future differently. Appendix \ref{appendix:rhovariance} sizes the alternative on the model with everything else common: reproducing one tenth of the variance of log default across score deciles would require a standard deviation of the discount rate across deciles of $0.13$ within twelve months and $0.18$ over the life, larger than the estimate itself, and no rate up to $3$ per year reproduces one half, whereas a factor of $2.6$ in the payment-to-income ratio, at the model's payment elasticity, reproduces the whole bottom-to-top tercile gap. Across the score distribution the same logic places the bottom deciles' defaults well past the threshold, the top deciles' among the few whose circumstances took them across it, and the middle deciles' near it, where the elasticity peaks. The heterogeneity in default is exposure, not impatience. Low-income borrowers default far more often because they enter into more burdensome contracts, the payment-to-income ratio in Table \ref{tab:summary_families} being highest where income is lowest, and face more volatile income and expense processes, not because they weigh the future less.
